\chapter{Graph-Augmented Vision Transformer}
\label{chap:method}

\section{Problem formulation and architectural scope}

This study investigates whether explicit region aggregation and graph reasoning
provide useful additional structure for remote sensing scene classification.
The comparison uses a common Swin-T backbone, so that the principal architectural
changes occur after the backbone feature extraction stage. Swin was proposed by
Liu et al.~\cite{liu2021swin}; its shifted-window hierarchy already supports
contextual interaction. The motivation here is therefore to test an additional
regional inductive bias, rather than to assume that the backbone lacks spatial
or contextual modelling.

Let $I\in\mathbb{R}^{3\times224\times224}$ denote a preprocessed RGB image, and
let $C$ denote the number of output classes. For single-label classification,
the target $y$ is one class index in $\{1,\ldots,C\}$. For multi-label
classification, the target is a binary vector
$\boldsymbol{y}\in\{0,1\}^{C}$. The backbone produces a sequence
$X=[x_1,\ldots,x_N]^{\mathsf T}\in\mathbb{R}^{N\times D}$, where $N=49$
and $D=768$. The tokens occupy a $7\times7$ spatial grid. They are contextual
features from the final backbone stage, rather than independent raw-image
patches.

We implement a Graph-Augmented Vision Transformer (GAViT) with four successive
operations: token-to-region aggregation, sparse region graph construction,
graph attention, and region-to-token feedback. The principal comparison uses
either a fixed spatial partition with learned pooling weights or a local,
image-dependent soft assignment. Both versions retain $K=16$ region anchors
and the same explicit graph-edge budget. The term \emph{region} denotes a group
of feature tokens; it does not imply a supervised object or land-cover segment.

\section{Shared feature extractor and baseline}

The backbone is the \texttt{swin\_tiny\_patch4\_window7\_224} model.
Its final feature map has spatial dimensions $7\times7$ and 768 channels and
is reshaped into $X$. The experiments fine-tune the backbone from the recorded
ImageNet-pretrained checkpoint. No additional classification token is inserted.

The baseline first computes the global representation
\begin{equation}
 \bar{x}=\frac{1}{N}\sum_{n=1}^{N}x_n,
 \label{eq:baseline-pooling}
\end{equation}
where $n$ indexes tokens. Its logit vector is
\begin{equation}
 \boldsymbol{\ell}=W_c\,\operatorname{Dropout}
       \!\left(\operatorname{LN}(\bar{x})\right)+b_c.
 \label{eq:classifier}
\end{equation}
Here $\operatorname{LN}$ denotes layer normalisation,
$W_c\in\mathbb{R}^{C\times D}$ and $b_c\in\mathbb{R}^{C}$ are classifier
parameters, and $\operatorname{Dropout}$ uses probability 0.1 during training
and is disabled during evaluation. GAViT uses the same form of classifier
after its feedback operation.

\section{Fixed spatial partition with attentive pooling}

\subsection{Deterministic token membership}

For the assignment formula only, region identifiers are zero-based. A token at
grid coordinate $(u,v)$, with $u,v\in\{0,\ldots,6\}$, is assigned to
\begin{equation}
 a(u,v)=4\left\lfloor\frac{4u}{7}\right\rfloor
          +\left\lfloor\frac{4v}{7}\right\rfloor.
 \label{eq:fixed-assignment}
\end{equation}
The function $a$ maps each token to one of the 16 anchors. We write $a(n)$ for
the assignment of token $n$. Along each grid axis, the resulting group sizes
are $[2,2,2,1]$. Thus, nine regions contain four tokens, six contain two tokens,
and one contains one token. This is a partition of a $7\times7$ feature grid;
it is not a division into 16 equally sized input-image crops.

\subsection{Learned pooling within each region}

A shared scorer produces one scalar $s_n=f_{\theta}(x_n)$ per token.
The scorer $f_{\theta}$ consists of a 768-to-192 linear layer, a Gaussian error
linear unit (GELU), and a 192-to-1 linear layer. Its trainable parameters are
denoted collectively by $\theta$. The fixed-membership pooling weights are
\begin{equation}
 w^{(0)}_{nk}=
 \frac{\mathbf{1}[a(n)=k]\exp(s_n)}
 {\sum_{m=1}^{N}\mathbf{1}[a(m)=k]\exp(s_m)},
 \qquad
 r_k=\sum_{n=1}^{N}w^{(0)}_{nk}x_n,
 \label{eq:fixed-pooling}
\end{equation}
where $k\in\{0,\ldots,K-1\}$ indexes regions, $m$ is a second token index,
$\mathbf{1}[\cdot]$ is an indicator, and
$r_k\in\mathbb{R}^{D}$ is the pooled region feature.
The weights sum to one over the tokens in each region. The scorer changes
which member tokens contribute most strongly, but it does not move the
partition boundaries or change the member set.

\section{Local prototype soft assignment}

The soft version changes token membership while preserving the anchor layout.
It constructs image-dependent prototypes from the fixed partition, measures
token--prototype similarity in a learned projection space, and restricts
assignment to nearby anchors. This is a single assignment step; it is not an
iterative clustering procedure or a bank of independently learned class
prototypes.

\subsection{Projected tokens and image-dependent prototypes}

Let $P\in\mathbb{R}^{128\times D}$ denote a learned linear projection without
bias. Define
\begin{equation}
 z_n=\operatorname{unit}\!\left(P\operatorname{LN}_{a}(x_n)\right),
 \qquad
 p_k=\operatorname{unit}\!\left(\sum_{n=1}^{N}w^{(0)}_{nk}z_n\right),
 \label{eq:soft-prototypes}
\end{equation}
where $\operatorname{LN}_{a}$ is the assignment projection's layer
normalisation, $z_n\in\mathbb{R}^{128}$ is a normalised token projection, and
$p_k\in\mathbb{R}^{128}$ is the current image's prototype for anchor $k$.
The operation $\operatorname{unit}(v)=v/\max(\|v\|_2,\epsilon)$ normalises
a vector $v$, with a small positive numerical stabiliser $\epsilon$.
The initial weights $w^{(0)}$ are those in Equation~\eqref{eq:fixed-pooling}.
Consequently, the prototypes change with the image and the backbone features.

\subsection{Local membership and separate aggregation weights}

Let $g(k)=(\lfloor k/4\rfloor,k\bmod4)$ be the grid coordinate of anchor $k$.
The candidate set for token $n$ is
\begin{equation}
 \mathcal{C}(n)=
 \{k:\|g(a(n))-g(k)\|_{\infty}\leq1\},
 \label{eq:local-candidates}
\end{equation}
where $\|\cdot\|_{\infty}$ measures the largest absolute coordinate difference.
This restriction permits assignment to adjacent anchors, including diagonal
neighbours, while maintaining a spatial prior. With temperature $\tau=0.1$,
the membership matrix $S\in\mathbb{R}^{N\times K}$ is
\begin{equation}
 S_{nk}=
 \begin{cases}
 \displaystyle\frac{\exp(z_n^{\mathsf T}p_k/\tau)}
 {\sum_{j\in\mathcal{C}(n)}\exp(z_n^{\mathsf T}p_j/\tau)},
       & k\in\mathcal{C}(n),\\[0.7em]
 0,    & k\notin\mathcal{C}(n).
 \end{cases}
 \label{eq:membership}
\end{equation}
Here $j$ indexes candidate regions, and $\sum_k S_{nk}=1$ for every token.

Membership and region pooling use different normalisations. Define the pooling
matrix $W\in\mathbb{R}^{N\times K}$ by
\begin{equation}
 W_{nk}=\frac{S_{nk}\exp(s_n)}
                {\sum_{m=1}^{N}S_{mk}\exp(s_m)},
 \qquad
 r_k=\sum_{n=1}^{N}W_{nk}x_n.
 \label{eq:soft-pooling}
\end{equation}
Thus, $S$ sums to one over regions for each token, whereas $W$ sums to one over
tokens for each region. The implementation computes the latter operation using
$\log S$ and a token-axis softmax. Region aggregation uses the original
768-dimensional tokens, rather than the 128-dimensional projection.

The projection and assignment normalisation add 99,840 trainable parameters
relative to the fixed grouping module. Gradients pass through the local
membership calculation, the image-dependent prototypes and the pooling
weights. The locality mask remains fixed. The resulting regions may overlap
in their token contributions, but they retain anchor identities and are not
unconstrained semantic segments.

\section{Sparse spatial and feature graph}

The graph is constructed once per image from its pooled region features.
Its node set is $\mathcal{V}=\{0,\ldots,K-1\}$. Let
$\mathcal{N}_{4}(i)$ denote the up, down, left and right anchor neighbours of
node $i$, restricted to the $4\times4$ grid. The directed spatial edge set is
\begin{equation}
 \mathcal{E}_{s}=\{(j\rightarrow i):j\in\mathcal{N}_{4}(i)\}.
 \label{eq:spatial-edges}
\end{equation}
It contains 48 edges. Each neighbouring anchor pair contributes both message
directions.

Feature edges provide a second source of connectivity. For every query node
$i$, the cosine score to candidate node $j$ is
\begin{equation}
 c_{ij}=\frac{r_i^{\mathsf T}r_j}
 {\max(\|r_i\|_2,\epsilon)\max(\|r_j\|_2,\epsilon)}.
 \label{eq:region-cosine}
\end{equation}
Let $\mathcal{T}_2(i)$ contain the two candidates with the highest scores
among $\mathcal{V}\setminus(\{i\}\cup\mathcal{N}_{4}(i))$. Equal scores
are resolved by preferring the smaller region index. The feature and combined
edge sets are
\begin{equation}
 \mathcal{E}_{f}=\{(j\rightarrow i):j\in\mathcal{T}_2(i)\},
 \qquad
 \mathcal{E}=\mathcal{E}_{s}\cup\mathcal{E}_{f}.
 \label{eq:feature-edges}
\end{equation}
There are 32 feature edges and 80 unique explicit non-self edges in total.
The count describes the graph supplied to the graph-attention operator,
before its internal edge processing. It is not a count of every message
handled internally by that operator.

An edge $j\rightarrow i$ supplies information from selected neighbour $j$ to
query $i$. The feature graph depends on the image, while the spatial graph
retains the fixed anchor adjacency in both grouping variants. In particular,
soft region geometry is not used to recompute spatial touching relationships
or region centroids. Cosine similarity selects connectivity; its numeric
value is not supplied as an additional message weight to the graph-attention
layers. The top-two selection is discrete, so differentiable soft membership
does not make the entire graph topology differentiable. Both graph-attention
layers share the graph constructed before message passing.

\section{Graph attention and token feedback}

Graph attention networks (GATs) were proposed by
Veli\v{c}kovi\'{c} et al.~\cite{velickovic2018gat} to learn neighbour
aggregation weights. Here, let $R\in\mathbb{R}^{K\times D}$ be the matrix
of region features. An input linear projection maps its rows to width 1024,
yielding $H^{(0)}\in\mathbb{R}^{K\times1024}$. Two graph-attention layers
then apply
\begin{equation}
 H^{(l+1)}=\operatorname{Dropout}\!\left(
  \operatorname{GELU}\!\left[
  \operatorname{LN}_{l}\!\left(
   \operatorname{GAT}_{l}(H^{(l)},\mathcal{E})+H^{(l)}
  \right)\right]\right),\quad l\in\{0,1\}.
 \label{eq:gat-block}
\end{equation}
The layer index is $l$, $\operatorname{LN}_{l}$ is its layer normalisation,
and $\operatorname{GAT}_{l}$ has four concatenated attention heads with 256
output channels per head. Both attention dropout and the output dropout use
probability 0.1. The residual connection preserves the 1024-dimensional shape.

Let $h_k^{(2)}$ denote row $k$ of $H^{(2)}$. A learned projection
$Q\in\mathbb{R}^{D\times1024}$ with bias $b_f\in\mathbb{R}^{D}$ and a
further layer normalisation $\operatorname{LN}_{f}$ produces
\begin{equation}
 F_k=\operatorname{LN}_{f}(Qh_k^{(2)}+b_f),
 \label{eq:region-feedback}
\end{equation}
where $F_k\in\mathbb{R}^{D}$ is the feedback for region $k$. The updated
tokens are
\begin{equation}
 x'_n=
 \begin{cases}
 x_n+F_{a(n)}, & \text{fixed grouping},\\
 x_n+\sum_{k=0}^{K-1}S_{nk}F_k, & \text{soft grouping}.
 \end{cases}
 \label{eq:token-feedback}
\end{equation}
No learned scalar gate is applied at this interface. The updated tokens are
mean-pooled and passed to the classifier in Equation~\eqref{eq:classifier}.
The principal model does not apply another Transformer block after feedback.

This placement also limits the interpretation of token feedback. In the soft
version, the vector passed to the classifier satisfies
\begin{equation}
 \frac{1}{N}\sum_{n=1}^{N}x'_n
 =\bar{x}+\sum_{k=0}^{K-1}\mu_kF_k,
 \qquad
 \mu_k=\frac{1}{N}\sum_{n=1}^{N}S_{nk}.
 \label{eq:pooled-feedback}
\end{equation}
Here $\mu_k$ is the average membership mass of region $k$. This identity
describes the final pooling step, not the complete effect of assignment:
$S$ also changes the upstream region features, feature-edge selection and
graph outputs. It nevertheless prevents an unsupported claim that feedback
is followed by further token-to-token reasoning in the principal model.

\section{Classification objectives}

For single-label tasks, let $q_c$ be the softmax probability of class $c$,
computed from logit $\ell_c$. The per-image cross-entropy objective is
\begin{equation}
 q_c=\frac{\exp(\ell_c)}{\sum_{j=1}^{C}\exp(\ell_j)},
 \qquad \mathcal{L}_{\mathrm{single}}=-\log q_y,
 \label{eq:single-loss}
\end{equation}
where $j$ now indexes output classes and $y$ is the target class.
For multi-label tasks, define $p_c=\sigma(\ell_c)$, where $\sigma$ is the
sigmoid function and $y_c$ is the binary target for class $c$. The objective is
\begin{equation}
 \mathcal{L}_{\mathrm{multi}}=-\frac{1}{C}\sum_{c=1}^{C}
 \left[y_c\log p_c+(1-y_c)\log(1-p_c)\right].
 \label{eq:multi-loss}
\end{equation}
Both objectives are averaged over the training minibatch. The implementation
uses binary cross-entropy directly on logits for numerical stability in the
multi-label case. No region masks, object annotations or auxiliary grouping
loss are used in the principal models. Consequently, classification training
alone does not establish that the learned grouping recovers semantic
boundaries. That question requires separate evidence from the classification
comparison and qualitative analysis.
